\begin{helper}[Fixed-fiber tail probability]
Fix \(n,m,p\), a red row or blue column indexed by \(x\in [m]\), and
\(\delta\in(0,1/2]\).  If \(n\le m^2\) and \(m>0\), then with
\(\mu=n/m\) the two-bite probability that this fixed fiber has size larger
than \((1+\delta)\mu\) is at most
\[
  \exp\!\left(-\bigl((1+\delta)\log(1+\delta)-\delta\bigr)\mu\right),
\]
and the two-bite probability that it has size smaller than
\((1-\delta)\mu\) is at most
\[
  \exp\!\left(-\bigl(\delta+(1-\delta)\log(1-\delta)\bigr)\mu\right).
\]
\end{helper}
\begin{proof}
For the chosen row or column let \(K\subseteq [m]\times[m]\) be the
corresponding marked set.  It has cardinality \(m\).  The node
\(\noderef{FiberAndDegreeFiberSizeInjectionMarginal}\) reduces the two-bite
probability of any fixed-fiber count event to the uniform injection marginal.
For each injection, the number of domain vertices in the fixed fiber is equal
to the size of the sampled image intersected with \(K\), because the injection
is one-to-one.  Thus the upper and lower count events are exactly the image
tail events for \(K\).  Applying
\(\noderef{FiberAndDegreeFiberSizeImageHypergeometricUpperTailChernoff}\) and
\(\noderef{FiberAndDegreeFiberSizeImageHypergeometricLowerTail}\), then
dividing by the total number of injections through
\(\noderef{UniformInjectionWeight}\), gives the two displayed bounds.
\end{proof}
